\chapter*{Contribuciones Personales}
\label{cap:contribucionesPersonales}
\addcontentsline{toc}{chapter}{Contribuciones Personales}

\section*{Eva Feliu Aréjola}
Dada la complejidad teórica del proyecto era esencial la coordinación de las distintas participantes, por lo que varias tareas se han llevado a cabo en colaboración, ya fuera de manera paralela o con intercambios consecutivos entre las colaboradoras del proyecto. Asímismo, la gestión logística del proyecto se ha desempeñado por todas las autoras: la comunicación con el tutor, la organización de los puntos a discutir en cada reunion, la redacción del acta de las mismas, la división de tareas y el seguimiento de estas en github, son algunos de los ejemplos. \newline

Eva comenzó junto a sus compañeras la investigación para acotar el tema de este proyecto, manteniendo un estudio sobre aplicaciones novedosas de la IA en físicas, y enfoques de simulaciones interesantes, particularmente en videojuegos. Una vez concretado el tema y contemplada la evolución general de los tipos y evolución de las simulaciones de telas, se centró en aquellas basadas en grandes datasets, sistemas de muelles, CNNs y reducciones de dimensionalidad. Por otro lado, junto a su compañera Paula, investigó las soluciones existentes al coste computacional más allá de la IA, viendo optimizaciones por CPU en herramientas o paralelización por GPU, centrándose en Unity como plaforma elegida. Se concluyó que herramientas no eran accesibles al público general, por complejidad de implentación, y por no ser gratuitos. \newline

Respecto a las simulaciones necesarias para entrenar la tela, realizó iteraciones adicionales sobre distintas mallas en Blender, para satisfacer requisitos concretos a lo largo del desarrollo, como reducciones o aumentos de dimensionalidad de la tela, y su posterior puesta en escena. \newline

Para poder guardar los datos obtenidos de la simulación, creó la base del \textit{ClothDatasetRecoder} junto a Paula, extrayendo los datos mencionados relevantes de la tela y guardándolos en formato csv. Se habilitó el grabado automático de segmentos de simulaciones, que generaba los consecuentes archivos. Eva investigó diferentes formas de tener en cuenta al colisionador en los datos como flags de contacto o campos voxelizados, y terminó creando un gran interés por los Signed Distance Fields, medidas de distancia a las superficies que permitían tener en cuenta un cuerpo formado por varios colisionadores y valorar volúmenes para simulaciones más complejas, como sería la falda. Para habilitar estos cálculos para el guardado de datos y posterior simulación con el modelo, implementó la clase estática \textit{SDFTUtil}, que permite realizar los cálculos de distancias para un conjunto de colliders de tipo esfera y cápsula, los aceptados por Cloth. Trabajó en traer el espacio de coordenadas al de la tela en todas estas funciones, después de las pruebas realizadas para elegir entre coordenadas globales y locales.\newline

Tras establecer Mikele la base de Pytorch con los datos y primer borrador de modelo, estableció el primer modelo MLP con las capas y tamaño básicos para comenzar la investigación. También habilitó la posibilidad de realizar el entrenamiento sobre un conjunto de datasets, pero esta característica sería descartada tras advertir la necesidad de secuencias temporales, más fácilmente alcanzables sobre una sola simulación larga. Eva confirmó el Adam como el optimizador adecuado para las simulaciones físicas.  Durante las primeras pruebas en el entorno de simulación, realizó los cambios necesarios para deformar la malla de \textit{ClothML} en Unity a tiempo real, revisar restricciones y conseguir un resultado visual. \newline

Sobre el MLP llevó a cabo distintas pruebas de tamaños, funciones de activación (cómo relu vs leakyrelu) así como una forma distinta de obtener los desplazamientos de los vértices. Esta última consiste en lo que se denomina como "Anchored", o en otras palabras, realizar todas las transformaciones de los vértices sobre la pose inicial en reposo de la malla, que propiciaba que el error acumulado disminuyera. Está linea de estudio terminó descartándose como modelo principal, al obtener resultados visualmente menos atractivos que la combinación de recurrencia y unfolding. En su lugar, aplicaría sobre el MLP la técnica de unfolding, para poder comparar en los experimentos el efecto de la recurrencia en rendimiento y efecto visual. \newline

Más avanzado el proyecto, otra parte del trabajo realizado por Eva consistió en llevar a cabo medidas de rendimiento, comparando inferencias entre distintos tipos de modelos y, particularmente, células recurrentes diferentes: RNN, GRU, LSTM. Lamentablemente, el paquete Sentis no admite células GRU, que tenían un rendimiento muy ligeramente mejor. En la misma línea del rendimiento, trabajó en la habilitación del entrenamiento por GPU y habilitó la captura de inferencia del Modelo en runtime en Unity.
Para poder observar el comportamiento de la tela en simulaciones de mayor coste, hizo pruebas con la simulación de la falda tomando diferentes complejidades (72, 288 vértices). Las pruebas mostraron que el modelo necesita un mucho mayor número de datos para entrenarse y mantener la topología, inaccesible con las limitaciones del proyecto. \newline

Eva creó la versión corregida del PCA, iniciado por Mikele, incorporando el PCA inverso: la extracción de posiciones locales mediante la transformación inversa en el subespacio, tanto en Unity como en el entorno de entrenamiento. Llevó a cabo pruebas con distintos tamaños de dataset, así como la adición de \textit{clipping} sobre los resultados a los obtenidos en el entrenamiento, con objeto de reducir movimientos erráticos producidos por datos no esperados. \newline

Y finalmente, como se ha mencionado como dinámica del equipo, se unió a tareas externas, resolución de bugs y modificaciones para conseguir llegar al punto final descrito en este proyecto, así como todos los experimentos necesarios para obtener las métricas de los experimentos y pruebas intermedias para el desarrollo. De la misma forma, colaboró activamente en la escritura de esta memoria, traducciones y sus posteriores revisiones.

\newpage
\section*{Paula López Román}
Dada la complejidad teórica del proyecto era esencial la coordinación de las distintas participantes, por lo que varias tareas se han llevado a cabo en colaboración, ya fuera de manera paralela o con intercambios consecutivos entre las colaboradoras del proyecto. Asímismo, la gestión logística del proyecto se ha desempeñado por todas las autoras: la comunicación con el tutor, la organización de los puntos a discutir en cada reunion, la redacción del acta de las mismas, la división de tareas y el seguimiento de estas en github, son algunos de los ejemplos. \newline

Para poder comenzar con la investigación del proyecto, Paula configuró una librería compartida en Zotero en la que se guardarían y organizarían todos los artículos leidos para asistir el desarrollo del proyecto. Con respecto al trabajo de investigación, es necesario recalcar que ha estado presente a lo largo de todo el desarrollo del proyecto, convirtiendolo en una tarea necesariamente compartida. En los primeros meses de desarrollo Paula, al igual que sus compañearas, investigó a cerca de las distintas investigaciones en curso relacionadas con el tema de estudio, hasta que se acabaron de concretar las necesidades específicas del proyecto. Posteriormente se indagó en mayor profundidad en las diversas investigaciones del estado del arte, en el que ella se especializó en las investigaciones dirigidas a PBD. \newline

Además de la investigación en el ámbito académico, Paula se encargó de hacer una búsqueda tanto teórica como práctica a través de diversos repositorios para entender la herramienta Sentis. Posteriormente introduciría esta herramienta al desarrollo, creando la base del respectivo componente en Unity. \newline

Antes de llegar a la fase de inferencia, la estudiante creó la base del movimiento del objeto colisionador en la escena de la simulación. Lo dotó de un movimiento lineal automático, para poder comprobar las primeras iteraciones del modelo con datasets sencillos con overfitting, y un movimiento generado por el input del usuario en los ejes X y Z. Las simulaciones seguirían siendo ajustadas indistintamente por todas las participantes del proyecto en los meses después de crear esta base; probando distintos parámetros para el componente de cloth de Unity, distintos tamaños o velocidades del colisionador, mallas con distintas topologías y demás adaptaciones.\newline

Para poder pasar a la fase de entrenamiento, fue esencial extraer los datos de cada vértice de la geometría a lo largo de los frames. Para ello se creó un sistema de grabación y registro, \textit{ClothDatasetRecoder}, basado en archivos CSV. Este trabajo fue llevado a cabo en estrecha colaboración con la compañera Eva Feliu Aréjola, entre las dos pudieron crear un sistema completo capaz de producir los datasets con la información pertinente. Partieron por diseñar la herramienta y las características necesarias para la grabación (que conformarían 13 features). Trabajó en un sistema de grabado por pasos en el que poder fijar el tiempo entre las capturas de la tela, asímismo, se optimizó el mecanismo implementando un volcado en memoria por intervalos. Uno de los aspectos sobre los que iteraron fue traer el espacio de coordenadas al de la tela en todas las funciones, después de las pruebas realizadas para elegir entre coordenadas globales y locales. En cuanto a la organización en el directorio destino, crearon un sistema incremental para el nombrado automático de las distintas grabaciones. \newline

Paula trabajó en la gestión de la normalización de los datos para su correcto entrenamiento. Para ello implementó la normalización tras la limpieza de los datasets, guardado las variables de normalización (la media y la desviación típica) en un archivo JSON para poder utilizar posteriormente en la fase de inferencia.  \newline 

En referencia al desarrollo de los diferentes modelos, Paula trabajo de manera paralela a sus compañeras con las diversas iteraciones. Particularmente encontro la técnica de injección de ruido a los datos de entrada del modelo, que resultó ser muy beneficiosa para el control de los errores. \newline

A su vez, y para permitir la ejecución de los modelos en un entorno real de simulación, Paula comenzó con la integración del ML a Unity. Como se ha mencionado previamente, dio los primeros pasos de la incorporación de Sentis al flujo de trabajo. De tal manera que consiguió reemplazar el componente de cloth propio de Unity con un componente nuevo, \textit{ClothML}. El componente gestionaba los anclajes de la geometría, la extracción de datos de entrada, su gestión a través de un \textit{history buffer} para respetar la arquitectura \textit{many to one} la normalización de los mismos a través del JSON previamente mencionado, la conversión a un tensor adecuado para el modelo cargado, la ejecución del modelo de sentis, la extracción de los deltas de desplazamiento y la actualización de la tela a través de ellos. Considerando la naturaleza del proyecto, se iteró sobre esta base para crear las adaptaciones a distintos modelos y diversas optimizaciones y pruebas. \newline

Una vez el proyecto contaba con unos resultados sólidos, Paula diseñó las pruebas de usuario que permitirían a gente externa al proyecto evaluar la calidad de la inferencia de los modelos más óptimos. Para ello generó distintos videos de las simulaciones, preparó un formulario en Google Forms y redactó el mensaje para su posterior distribución. \newline

En las últimas semanas de desarrollo colaboró con sus compañeras en llevar a cabo los experimentos diseñados y analizar los resultados. Así como la finalización de las escenas de los experimentos para la aplicación. Finalmente, como se ha mencionado como dinámica del equipo, se unió a tareas externas, resolución de bugs y modificaciones para conseguir llegar al punto final descrito en este proyecto, así como todos los experimentos necesarios para obtener las métricas de los experimentos y pruebas intermedias para el desarrollo. De la misma forma, colaboró activamente en la escritura de esta memoria, traducciones y sus posteriores revisiones.

\newpage
\section*{Mikele López de la Hoz}
Dada la complejidad teórica del proyecto era esencial la coordinación de las distintas participantes, por lo que varias tareas se han llevado a cabo en colaboración, ya fuera de manera paralela o con intercambios consecutivos entre las colaboradoras del proyecto. Asímismo, la gestión logística del proyecto se ha desempeñado por todas las autoras: la comunicación con el tutor, la organización de los puntos a discutir en cada reunion, la redacción del acta de las mismas, la división de tareas y el seguimiento de estas en github, son algunos de los ejemplos. \newline

Al inicio del proyecto se hizo un esfuerzo colectivo por investigar a cerca de las distintas investigaciones en curso relacionadas con el tema de estudio, hasta que se acabaron de concretar las necesidades específicas del proyecto. Posteriormente se indagó en mayor profundidad en las diversas investigaciones del estado del arte, en particular, Mikele se centró en las más recientes, como SNUG o HOOD, para intentar determinar de manera más informada el alcance de la investigación y las técnicas aplicables. \newline

Implementó sus conocimientos en modelado procedural para conseguir rápida y flexiblemente algunos de los modelos necesarios para la generación de las simulaciones, estos modelos fueron creados utilizando el software de Houdini y posteriormente exportados a Unity. Otra de las tareas necesarias para la configuración de la escena de Unity fue aportar de aleatoriedad a la simulación. Para ello, ella implementó un sistema de generación de splines consecutivos que el objeto colisionador utilizaría como carril para sus movimientos. Las simulaciones seguirían siendo ajustadas indistintamente por todas las participantes del proyecto en los meses después de crear esta base; probando distintos parámetros para el componente de cloth de Unity, distintos tamaños o velocidades del colisionador, mallas con distintas topologías y demás adaptaciones.  \newline

A la par que sus compañeras avanzaban con el desarrollo de la herramienta de grabación de las simulaciones, Mikele compenzó a indagar en la librería Pytorch. Al ser un entorno de trabajo nuevo para las integrantes, fue importante conseguir recursos de aprendizaje fiables para implementar la estructura base del entrenamiento. Uno de los aspectos más desafiantes a los tuvo que hacer frente fue la creación del environment de Conda, ya que tuvo que sincronizar las distintas librerías necesarias para el desarrollo de los modelos. Una vez conseguido un entorno de trabajo favorable, creó la primera iteración de los pasos del entrenamiento: una limpieza básica de los datos, un modelo simple que recibiera un dataset y lo entrenara y una serie de mátricas para evaluar estos modelos. Partiendo de esta base todas contribuyeron al desarrollo de estos scripts para implementar cada modelo experimental. \newline

En cuanto a la evaluación de los modelos para hacer frente a los diversos errores encontrados en el proceso de desarrollo, Mikele empezó por implementar un modelo RNN. Este abrió un camino hacia las técnicas de sensibilización temporal. Llevó a cabo varias iteraciones para refinar la inferencia, sin embargo, el entrenamiento no conseguía aprender las dependencias temporales de manera óptima. Fue entonces que encontró la tecnica de unrolling, marcando un antes y después en el estudio. Esta técnica doto al modelo de la capacidad de autocorregirse. Similarmente, tuvieron que hacerse varias iteraciones para conseguir el modelo final, en lo que las compañeras participaron activamente. De la misma manera, en los pasos más avanzados del desarrollo, se hizo un esfuerzo congunto por valorar y elegir diferentes técnicas implementadas, las diferentes funciones de pérdida son un ejemplo de ello. En estos casos en los que más allá de la implementación hizo falta discutir los beneficios o perjuicios de distintos métodos, se participó en las pruebas y deliberación de manera colaborativa, llegando así a los modelos finales. \newline

Mikele trabajó en varias de las técnicas fallidas implementadas, parte indispensable de un proyecto tan experimental como este. Probó distintas variables que podrían ayudar al entrenamiento, implemento un modelo CNN y adaptó la entrada de datos a este y llevó a cabo el flujo de trabajo completo (preparación de la simulación, grabación de datos, entrenamiento del modelo, incorporación en Unity y análisis) para una simulación estática controlada por 4 puntos con movimiento aleatorio.  \newline

Por último en lo que se refiere al desarrollo inicial del entrenamiento también implementó las técnicas de elección de parámetros a través de modelos Random Forest. Posteriormente, se unió a sus compañeras en el ajuste de hiperparámetros y perfeccionamiento de los modelos. \newline

En cuanto a la fase de indiferencia, Mikele adaptó la base proporcionada por su compañera Paula para adecuarla a la naturaleza del proyecto. Por un lado, dividió el componente en una base común a la inferencia de todos los modelos y distintos scripts para cada modelo con los métodos pertinentes. Por otro lado, optimizó la generación del tensor de input con una mejor gestión de la memoria dinámica del programa.  \newline

En los últimos meses de desarrollo, Mikele empezó con la implementación de un sistema de recogida de métricas de rendimiento en una escena ideada para este tipo de pruebas. Posteriormente la convertiría en una aplicación completa con una interfaz para ejecutar los distintos experimentos del estudio. Eva colaboró con la implementación de la recogida del tiempo de inferencia de cada modelo. Consiguiendo así un apoyo visual e interactivo al resultado teórico del proyecto. \newline

En las últimas semanas de desarrollo Mikele colaboró con sus compañeras en llevar a cabo los experimentos diseñados y analizar los resultados. También genero los gráficos para el análisis de las pruebas de usuario, y su posterior discusión. De la misma forma, colaboró activamente en la escritura de esta memoria, traducciones y sus posteriores revisiones.